%% sections/h-sponsor-cosponsor.tex - H. R. 9510 Bill v4.0 (full bill) - Appendix H
%% Genre: recruitment / sign-on packet (checklist + signature lines). Source deliverable 08.
\billsec{APPENDIX H.}{SPONSOR AND COSPONSOR PACKET (NON-OPERATIVE).}\phantomsection\label{toc:aH}

\uscprov{1}{\textit{Independent research draft. Not pending legislation and not legal
advice; not endorsed by the FDA, HHS, or any Member of Congress. No Member has agreed
to sponsor this measure. Bill: H. R. 9510 (illustrative number; the Clerk assigns the
real number at introduction), 119th Congress, 2d Session. Refer to: House Committee on
Energy and Commerce (Subcommittee on Health). Governing rule: Rule XII of the Rules of
the House of Representatives.}}

\uscprov{1}{\textbf{Purpose of this packet.} This packet is the recruitment and
sign-on instrument for the draft bill. In one place it tells a prospective sponsoring
office what role the bill needs filled, what each signer is being asked to do, where
the measure is expected to be referred, and what must be in hand before introduction.
It does not introduce anything. Only a sitting Member of the House can introduce a
House bill; an outside research project cannot file it directly. H. R. 9510 is
illustrative, and the objective is singular: place the finished text in the hands of
one Member who will introduce it and refer it to the Committee on Energy and Commerce,
and gather original cosponsors to accompany it \cite{house-rule-xii,
congress-howlawsmade}.}

\uscprov{1}{\textbf{The sponsor.} A House bill has exactly one sponsor: the Member who
introduces it and whose name leads the bill. Every other supporter is a cosponsor. The
sponsor's office owns the introduced text and is the contact for the Clerk, the
Parliamentarian, and the committee of referral. Two items are required from the
sponsor before introduction. First, the sponsor's signature on the bill, which under
current cloakroom guidance is required on the measure itself. Second, a signed
Constitutional Authority Statement: under Rule XII, clause 7(c)(1), a bill may not be
introduced unless the sponsor submits a statement citing, as specifically as
practicable, the constitutional power to enact it, here Article I, Section 8, Clause 3
(the Commerce Clause) on which the Act rests, supported by Clause 18 (the Necessary
and Proper Clause). The best-fit sponsor is a Member of the Committee on Energy and
Commerce, ideally on the Subcommittee on Health, with a demonstrated interest in
medical-device safety, oncology care, or trustworthy AI; that office is also expected
to route the final text through the House Office of the Legislative Counsel for proper
form before filing \cite{house-rule-xii, house-holc}.}

\uscprov{1}{\textbf{Original cosponsors.} Original cosponsors are Members listed at the
time of introduction, distinct from cosponsors added afterward. A public bill may have
an unlimited number of cosponsors, and under Rule XII the sponsor names them.
Bipartisan original cosponsors meaningfully improve a bill's prospects, so recruiting
at least one Member from each party is a deliberate goal. To be counted as original, a
Member must commit before or at introduction; once the bill is filed the
original-cosponsor list closes and later supporters are added through the eHopper.
Circulate for sign-on only after the sponsor approves the exact introduction version,
so every signer endorses identical text. A balanced slate spans three interests on
both sides of the aisle: device and FDA oversight, oncology and patient safety, and AI
and emerging technology \cite{house-rule-xii, house-ehopper}.}

\uscprov{1}{\textbf{Committee referral.} Expect primary referral to the Committee on
Energy and Commerce (Subcommittee on Health), which holds jurisdiction over the
Federal Food, Drug, and Cosmetic Act and FDA device authorities. Because the bill
addresses machine-generated robot-patient interaction code and the autonomy of
Physical AI systems, a secondary referral to the Committee on Science, Space, and
Technology is possible. The Speaker, advised by the Parliamentarian, makes the
referral after introduction; the sponsor's staff may ask the Parliamentarian to review
the draft in advance for likely referral and parliamentary problems.}

\uscprov{1}{\textbf{Introduction checklist.} Each item is confirmed before the bill is
filed.}

\cboxx\ One authoritative introduction version of the text is selected, frozen, and
version-controlled, with no further edits to the introduced version.

\cboxx\ A named human attorney and qualified oncology and regulatory experts have
reviewed and taken ownership of the final text.

\cboxx\ Currency check completed: no newer public law has altered Title 21, the target
section number (515D, 21 U.S.C. 360e-5) remains unused, and every cross-reference and
conforming amendment still lands.

\cboxx\ Final text routed through the House Office of the Legislative Counsel for
proper legislative form.

\cboxx\ Parliamentarian review requested, through the sponsor's staff, for likely
referral and parliamentary issues.

\cboxx\ Sponsor identified and committed, and the sponsor's signature obtained on the
bill.

\cboxx\ Constitutional Authority Statement finalized, signed, and ready to submit with
the measure (Article I, Section 8, Clause 3, supported by Clause 18).

\cboxx\ Original cosponsor list collected on the exact introduction version, with
bipartisan balance pursued.

\cboxx\ Companion materials assembled: one-pager, section-by-section, policy memo,
findings, and Ramseyer print.

\cboxx\ PAYGO statement and earmark declaration prepared as applicable; the bill
directs rulemaking only and is expected to be PAYGO neutral.

\cboxx\ Filing logistics confirmed: submission via the eHopper, the House's primary
electronic introduction mechanism, accepted from about 15 minutes before convening
through about 15 minutes after adjournment.

\cboxx\ AI-assisted origin disclosed to the sponsoring office, with a named human owner
taking responsibility for the final text.

\uscprov{1}{\textbf{Sign-on.} The first signer is the sponsor (one only). All
subsequent signers committing before or at introduction are original cosponsors. These
lines are illustrative placeholders; no Member has agreed to sign, and only a sitting
Member may introduce the measure.}

\sigline{Sponsor (Representative), date}

\sigline{Original cosponsor (Representative), date}

\sigline{Original cosponsor (Representative), date}

\uscprov{1}{\textbf{Transparency note.} This draft and its companion materials were
AI-assisted, generated with Claude Code (Opus 4.8). No House rule requires disclosure
that text was AI-drafted, but the sponsoring office should be told plainly, and a
named human attorney and subject-matter experts should review the text and take
ownership of it before introduction. Until a Member actually introduces the bill, it
remains an independent research draft: not pending legislation, not endorsed by the
FDA, HHS, or any Member of Congress, and not legal advice \cite{anthropic-claude-code}.}
